% 4.电磁学.tex 子文件，book文档类，主文档通过% 4.电磁学.tex 子文件，book文档类，主文档通过% 4.电磁学.tex 子文件，book文档类，主文档通过% 4.电磁学.tex 子文件，book文档类，主文档通过\input{4.电磁学.tex}导入
\part{第二篇 电磁学}
\setcounter{chapter}{5}
\chapter{第六章 电荷与电场}

\section{库仑定律与电场强度}
    \subsection{电荷及其性质}
    \subsection{库仑定律}
    \[
    F=k\frac{q_1q_2}{r^3}\boldsymbol{r}
    \]
    比例系数$k$写成
    \[
    k=\frac{1}{4\pi\varepsilon_0}
    \]
    式中，$\varepsilon_0$称为真空介电常量，也叫真空电容率。
    %\[
    %\varepsilon_0\approx 8.85\times10^{-12}\,\mathrm{F/m}
    %\]
    \subsection{电场与电场强度}
    \subsection{场强叠加原理}

\section{电场强度通量与高斯定理}
    \subsection{电场线}
    \subsection{电场强度通量}
    通过某个闭合曲面$S$的$\boldsymbol{E}$通量
    \subsection{真空中高斯定理}
    在真空中静电场内通过任意闭合曲面的$\boldsymbol{E}$通量，等于该闭合曲面所包围的电荷的代数和除以，其数学表达式为
    %\[
    %\]
    若场源电荷连续分布，则式为
    %\[
    %\]
    \subsection{高斯定理的应用}

\section{静电场的环路定理与电势}
    \subsection{静电场的环路定理}
    在静电场中，场强沿闭合路径的线积分等于零，称静电场环路定理。
    \subsection{电势}
    \subsection{电势差}
    \subsection{电势的计算}
    \subsection{等势面 电势梯度}
        \subsubsection{等势面}
        \subsubsection{电势梯度}

\section{静电场中导体}
    \subsection{导体的静电平衡性质}
        \subsubsection{导体的静电平衡状态}
        \subsubsection{静电平衡下导体带电特性}
    \subsection{空腔导体和静电屏蔽}
        \subsubsection{空腔导体}
        \subsubsection{静电屏蔽}
    \subsection{电容和电容器}
        \subsubsection{孤立导体的电容}
        \subsubsection{电容器及其电容}
        \subsubsection{串联和并联}
            \subsubsubsection{串联}
            \subsubsubsection{并联}

\section{静电场中的电介质 电介质中的高斯定理}
    \subsection{电介质的微观结构和极化过程}
    \subsection{电极化强度矢量}
    \subsection{电介质中的场强}
    \subsection{电介质中的高斯定理 电位移$\boldsymbol{D}$}
    \subsection{有电介质时静电场的计算}

\section{电场的能量}
    \subsection{点电荷系统的电能}
    \subsection{电容器的能量}
    \subsection{电场能量 电场能量密度}
    \subsection{电场能量的计算}
    电场能量为
    %\[
    %d
    %\]
    即可求得整个电场的总能量为
    %\[
    %W_e
    %\]
    显然

\chapter{第七章 电流与磁场}

\section{恒定电流和恒定电场 电动势}
    \subsection{电流形成的条件}
    定义单位时间内通过导线任一截面的电荷量为电流。

    在导体中形成电流的条件是：（1）导体（1）
    \subsection{恒定电流和恒定电场}
    \subsection{电流和电流密度}
    定义单位时间内通过垂直于电流方向单位面积的电量为导体中某点电流密度矢量的$\boldsymbol{j}$大小，$\boldsymbol{j}$的方向与正电荷在该点
    \subsection{欧姆定律的微分形式 焦耳—楞次定律的微分形式}
    \subsection{电源及电源电动势}
    电源电动势在量值

\section{恒定磁场和磁感应强度}
    \subsection{磁性起源于电荷的运动}
    \subsection{磁场 磁感应强度$\boldsymbol{B}$}

\section{毕奥‑萨伐尔定律}
    \subsection{毕奥‑萨伐尔定律}
    在真空中
    \subsection{毕奥‑萨伐尔定律应用举例}
    \subsection{匀速运动电荷的磁场}

\section{磁场中的高斯定理}
    \subsection{磁感应线}
    \subsection{磁通量}
    \subsection{恒定磁场中的高斯定理}

\section{真空中恒定磁场的安培环路定理}
    \subsection{恒定磁场的安培环路定理}
    \subsection{安培环路定理的应用}

\section{磁场对运动电荷和载流导线的作用}
    \subsection{洛伦兹力}
    \subsection{带电粒子在磁场中的运动}
    \subsection{应用电场和磁场控制带电粒子运动的实例}
    \subsection{安培力}

\section{磁力的功}
    \subsection{磁力对运动载流导线的功}
    \subsection{磁力矩对转动载流线圈的功}
    磁力矩对转动载流线圈的功等于电流强度与回路包围面积内的磁通量的增量的乘积。

\section{磁介质中的恒定磁场}
    \subsection{磁介质及其磁化}
    \subsection{磁介质中的高斯定理和安培环路定理}
    \subsection{铁磁质}

\chapter{第八章 电磁场与麦克斯韦电磁场方程组}

\section{电磁感应的基本定律}
    当穿过闭合导体回路的磁通量发生变化时，闭合导体回路中就会出现电流，称之为电磁感应。

    所以

    导体回路中

    应该

    如果，则
    %\[
    %\]
    在$t_1$到$t_2$这段时间为
    %\[
    %\]
    式表明

    总电动势为
    %\[
    %\]

\section{动生电动势}
    把感应电动势
    \subsection{动生电动势的产生原因}
    \subsection{动生电动势的计算}

\section{感生电动势和感生电场}
    \subsection{感生电动势}
    \subsection{感生电场}
    \subsection{感生电动势的计算}
    \subsection{涡电流}
        \subsubsection{热效应}
        \subsubsection{机械效应}
        \subsubsection{趋肤效应}

\section{自感和互感}
    \subsection{自感}
    \subsection{互感}

\section{磁场的能量}

\section{位移电流和全电流定律}
    \subsection{位移电流}
    \subsection{全电流定律}

\section{麦克斯韦电磁场方程组}

\endinput导入
\part{第二篇 电磁学}
\setcounter{chapter}{5}
\chapter{第六章 电荷与电场}

\section{库仑定律与电场强度}
    \subsection{电荷及其性质}
    \subsection{库仑定律}
    \[
    F=k\frac{q_1q_2}{r^3}\boldsymbol{r}
    \]
    比例系数$k$写成
    \[
    k=\frac{1}{4\pi\varepsilon_0}
    \]
    式中，$\varepsilon_0$称为真空介电常量，也叫真空电容率。
    %\[
    %\varepsilon_0\approx 8.85\times10^{-12}\,\mathrm{F/m}
    %\]
    \subsection{电场与电场强度}
    \subsection{场强叠加原理}

\section{电场强度通量与高斯定理}
    \subsection{电场线}
    \subsection{电场强度通量}
    通过某个闭合曲面$S$的$\boldsymbol{E}$通量
    \subsection{真空中高斯定理}
    在真空中静电场内通过任意闭合曲面的$\boldsymbol{E}$通量，等于该闭合曲面所包围的电荷的代数和除以，其数学表达式为
    %\[
    %\]
    若场源电荷连续分布，则式为
    %\[
    %\]
    \subsection{高斯定理的应用}

\section{静电场的环路定理与电势}
    \subsection{静电场的环路定理}
    在静电场中，场强沿闭合路径的线积分等于零，称静电场环路定理。
    \subsection{电势}
    \subsection{电势差}
    \subsection{电势的计算}
    \subsection{等势面 电势梯度}
        \subsubsection{等势面}
        \subsubsection{电势梯度}

\section{静电场中导体}
    \subsection{导体的静电平衡性质}
        \subsubsection{导体的静电平衡状态}
        \subsubsection{静电平衡下导体带电特性}
    \subsection{空腔导体和静电屏蔽}
        \subsubsection{空腔导体}
        \subsubsection{静电屏蔽}
    \subsection{电容和电容器}
        \subsubsection{孤立导体的电容}
        \subsubsection{电容器及其电容}
        \subsubsection{串联和并联}
            \subsubsubsection{串联}
            \subsubsubsection{并联}

\section{静电场中的电介质 电介质中的高斯定理}
    \subsection{电介质的微观结构和极化过程}
    \subsection{电极化强度矢量}
    \subsection{电介质中的场强}
    \subsection{电介质中的高斯定理 电位移$\boldsymbol{D}$}
    \subsection{有电介质时静电场的计算}

\section{电场的能量}
    \subsection{点电荷系统的电能}
    \subsection{电容器的能量}
    \subsection{电场能量 电场能量密度}
    \subsection{电场能量的计算}
    电场能量为
    %\[
    %d
    %\]
    即可求得整个电场的总能量为
    %\[
    %W_e
    %\]
    显然

\chapter{第七章 电流与磁场}

\section{恒定电流和恒定电场 电动势}
    \subsection{电流形成的条件}
    定义单位时间内通过导线任一截面的电荷量为电流。

    在导体中形成电流的条件是：（1）导体（1）
    \subsection{恒定电流和恒定电场}
    \subsection{电流和电流密度}
    定义单位时间内通过垂直于电流方向单位面积的电量为导体中某点电流密度矢量的$\boldsymbol{j}$大小，$\boldsymbol{j}$的方向与正电荷在该点
    \subsection{欧姆定律的微分形式 焦耳—楞次定律的微分形式}
    \subsection{电源及电源电动势}
    电源电动势在量值

\section{恒定磁场和磁感应强度}
    \subsection{磁性起源于电荷的运动}
    \subsection{磁场 磁感应强度$\boldsymbol{B}$}

\section{毕奥‑萨伐尔定律}
    \subsection{毕奥‑萨伐尔定律}
    在真空中
    \subsection{毕奥‑萨伐尔定律应用举例}
    \subsection{匀速运动电荷的磁场}

\section{磁场中的高斯定理}
    \subsection{磁感应线}
    \subsection{磁通量}
    \subsection{恒定磁场中的高斯定理}

\section{真空中恒定磁场的安培环路定理}
    \subsection{恒定磁场的安培环路定理}
    \subsection{安培环路定理的应用}

\section{磁场对运动电荷和载流导线的作用}
    \subsection{洛伦兹力}
    \subsection{带电粒子在磁场中的运动}
    \subsection{应用电场和磁场控制带电粒子运动的实例}
    \subsection{安培力}

\section{磁力的功}
    \subsection{磁力对运动载流导线的功}
    \subsection{磁力矩对转动载流线圈的功}
    磁力矩对转动载流线圈的功等于电流强度与回路包围面积内的磁通量的增量的乘积。

\section{磁介质中的恒定磁场}
    \subsection{磁介质及其磁化}
    \subsection{磁介质中的高斯定理和安培环路定理}
    \subsection{铁磁质}

\chapter{第八章 电磁场与麦克斯韦电磁场方程组}

\section{电磁感应的基本定律}
    当穿过闭合导体回路的磁通量发生变化时，闭合导体回路中就会出现电流，称之为电磁感应。

    所以

    导体回路中

    应该

    如果，则
    %\[
    %\]
    在$t_1$到$t_2$这段时间为
    %\[
    %\]
    式表明

    总电动势为
    %\[
    %\]

\section{动生电动势}
    把感应电动势
    \subsection{动生电动势的产生原因}
    \subsection{动生电动势的计算}

\section{感生电动势和感生电场}
    \subsection{感生电动势}
    \subsection{感生电场}
    \subsection{感生电动势的计算}
    \subsection{涡电流}
        \subsubsection{热效应}
        \subsubsection{机械效应}
        \subsubsection{趋肤效应}

\section{自感和互感}
    \subsection{自感}
    \subsection{互感}

\section{磁场的能量}

\section{位移电流和全电流定律}
    \subsection{位移电流}
    \subsection{全电流定律}

\section{麦克斯韦电磁场方程组}

\endinput导入
\part{第二篇 电磁学}
\setcounter{chapter}{5}
\chapter{第六章 电荷与电场}

\section{库仑定律与电场强度}
    \subsection{电荷及其性质}
    \subsection{库仑定律}
    \[
    F=k\frac{q_1q_2}{r^3}\boldsymbol{r}
    \]
    比例系数$k$写成
    \[
    k=\frac{1}{4\pi\varepsilon_0}
    \]
    式中，$\varepsilon_0$称为真空介电常量，也叫真空电容率。
    %\[
    %\varepsilon_0\approx 8.85\times10^{-12}\,\mathrm{F/m}
    %\]
    \subsection{电场与电场强度}
    \subsection{场强叠加原理}

\section{电场强度通量与高斯定理}
    \subsection{电场线}
    \subsection{电场强度通量}
    通过某个闭合曲面$S$的$\boldsymbol{E}$通量
    \subsection{真空中高斯定理}
    在真空中静电场内通过任意闭合曲面的$\boldsymbol{E}$通量，等于该闭合曲面所包围的电荷的代数和除以，其数学表达式为
    %\[
    %\]
    若场源电荷连续分布，则式为
    %\[
    %\]
    \subsection{高斯定理的应用}

\section{静电场的环路定理与电势}
    \subsection{静电场的环路定理}
    在静电场中，场强沿闭合路径的线积分等于零，称静电场环路定理。
    \subsection{电势}
    \subsection{电势差}
    \subsection{电势的计算}
    \subsection{等势面 电势梯度}
        \subsubsection{等势面}
        \subsubsection{电势梯度}

\section{静电场中导体}
    \subsection{导体的静电平衡性质}
        \subsubsection{导体的静电平衡状态}
        \subsubsection{静电平衡下导体带电特性}
    \subsection{空腔导体和静电屏蔽}
        \subsubsection{空腔导体}
        \subsubsection{静电屏蔽}
    \subsection{电容和电容器}
        \subsubsection{孤立导体的电容}
        \subsubsection{电容器及其电容}
        \subsubsection{串联和并联}
            \subsubsubsection{串联}
            \subsubsubsection{并联}

\section{静电场中的电介质 电介质中的高斯定理}
    \subsection{电介质的微观结构和极化过程}
    \subsection{电极化强度矢量}
    \subsection{电介质中的场强}
    \subsection{电介质中的高斯定理 电位移$\boldsymbol{D}$}
    \subsection{有电介质时静电场的计算}

\section{电场的能量}
    \subsection{点电荷系统的电能}
    \subsection{电容器的能量}
    \subsection{电场能量 电场能量密度}
    \subsection{电场能量的计算}
    电场能量为
    %\[
    %d
    %\]
    即可求得整个电场的总能量为
    %\[
    %W_e
    %\]
    显然

\chapter{第七章 电流与磁场}

\section{恒定电流和恒定电场 电动势}
    \subsection{电流形成的条件}
    定义单位时间内通过导线任一截面的电荷量为电流。

    在导体中形成电流的条件是：（1）导体（1）
    \subsection{恒定电流和恒定电场}
    \subsection{电流和电流密度}
    定义单位时间内通过垂直于电流方向单位面积的电量为导体中某点电流密度矢量的$\boldsymbol{j}$大小，$\boldsymbol{j}$的方向与正电荷在该点
    \subsection{欧姆定律的微分形式 焦耳—楞次定律的微分形式}
    \subsection{电源及电源电动势}
    电源电动势在量值

\section{恒定磁场和磁感应强度}
    \subsection{磁性起源于电荷的运动}
    \subsection{磁场 磁感应强度$\boldsymbol{B}$}

\section{毕奥‑萨伐尔定律}
    \subsection{毕奥‑萨伐尔定律}
    在真空中
    \subsection{毕奥‑萨伐尔定律应用举例}
    \subsection{匀速运动电荷的磁场}

\section{磁场中的高斯定理}
    \subsection{磁感应线}
    \subsection{磁通量}
    \subsection{恒定磁场中的高斯定理}

\section{真空中恒定磁场的安培环路定理}
    \subsection{恒定磁场的安培环路定理}
    \subsection{安培环路定理的应用}

\section{磁场对运动电荷和载流导线的作用}
    \subsection{洛伦兹力}
    \subsection{带电粒子在磁场中的运动}
    \subsection{应用电场和磁场控制带电粒子运动的实例}
    \subsection{安培力}

\section{磁力的功}
    \subsection{磁力对运动载流导线的功}
    \subsection{磁力矩对转动载流线圈的功}
    磁力矩对转动载流线圈的功等于电流强度与回路包围面积内的磁通量的增量的乘积。

\section{磁介质中的恒定磁场}
    \subsection{磁介质及其磁化}
    \subsection{磁介质中的高斯定理和安培环路定理}
    \subsection{铁磁质}

\chapter{第八章 电磁场与麦克斯韦电磁场方程组}

\section{电磁感应的基本定律}
    当穿过闭合导体回路的磁通量发生变化时，闭合导体回路中就会出现电流，称之为电磁感应。

    所以

    导体回路中

    应该

    如果，则
    %\[
    %\]
    在$t_1$到$t_2$这段时间为
    %\[
    %\]
    式表明

    总电动势为
    %\[
    %\]

\section{动生电动势}
    把感应电动势
    \subsection{动生电动势的产生原因}
    \subsection{动生电动势的计算}

\section{感生电动势和感生电场}
    \subsection{感生电动势}
    \subsection{感生电场}
    \subsection{感生电动势的计算}
    \subsection{涡电流}
        \subsubsection{热效应}
        \subsubsection{机械效应}
        \subsubsection{趋肤效应}

\section{自感和互感}
    \subsection{自感}
    \subsection{互感}

\section{磁场的能量}

\section{位移电流和全电流定律}
    \subsection{位移电流}
    \subsection{全电流定律}

\section{麦克斯韦电磁场方程组}

\endinput导入
\part{第二篇 电磁学}
\setcounter{chapter}{5}
\chapter{第六章 电荷与电场}

\section{库仑定律与电场强度}
    \subsection{电荷及其性质}
    \subsection{库仑定律}
    \[
    F=k\frac{q_1q_2}{r^3}\boldsymbol{r}
    \]
    比例系数$k$写成
    \[
    k=\frac{1}{4\pi\varepsilon_0}
    \]
    式中，$\varepsilon_0$称为真空介电常量，也叫真空电容率。
    %\[
    %\varepsilon_0\approx 8.85\times10^{-12}\,\mathrm{F/m}
    %\]
    \subsection{电场与电场强度}
    \subsection{场强叠加原理}

\section{电场强度通量与高斯定理}
    \subsection{电场线}
    \subsection{电场强度通量}
    通过某个闭合曲面$S$的$\boldsymbol{E}$通量
    \subsection{真空中高斯定理}
    在真空中静电场内通过任意闭合曲面的$\boldsymbol{E}$通量，等于该闭合曲面所包围的电荷的代数和除以，其数学表达式为
    %\[
    %\]
    若场源电荷连续分布，则式为
    %\[
    %\]
    \subsection{高斯定理的应用}

\section{静电场的环路定理与电势}
    \subsection{静电场的环路定理}
    在静电场中，场强沿闭合路径的线积分等于零，称静电场环路定理。
    \subsection{电势}
    \subsection{电势差}
    \subsection{电势的计算}
    \subsection{等势面 电势梯度}
        \subsubsection{等势面}
        \subsubsection{电势梯度}

\section{静电场中导体}
    \subsection{导体的静电平衡性质}
        \subsubsection{导体的静电平衡状态}
        \subsubsection{静电平衡下导体带电特性}
    \subsection{空腔导体和静电屏蔽}
        \subsubsection{空腔导体}
        \subsubsection{静电屏蔽}
    \subsection{电容和电容器}
        \subsubsection{孤立导体的电容}
        \subsubsection{电容器及其电容}
        \subsubsection{串联和并联}
            \subsubsubsection{串联}
            \subsubsubsection{并联}

\section{静电场中的电介质 电介质中的高斯定理}
    \subsection{电介质的微观结构和极化过程}
    \subsection{电极化强度矢量}
    \subsection{电介质中的场强}
    \subsection{电介质中的高斯定理 电位移$\boldsymbol{D}$}
    \subsection{有电介质时静电场的计算}

\section{电场的能量}
    \subsection{点电荷系统的电能}
    \subsection{电容器的能量}
    \subsection{电场能量 电场能量密度}
    \subsection{电场能量的计算}
    电场能量为
    %\[
    %d
    %\]
    即可求得整个电场的总能量为
    %\[
    %W_e
    %\]
    显然

\chapter{第七章 电流与磁场}

\section{恒定电流和恒定电场 电动势}
    \subsection{电流形成的条件}
    定义单位时间内通过导线任一截面的电荷量为电流。

    在导体中形成电流的条件是：（1）导体（1）
    \subsection{恒定电流和恒定电场}
    \subsection{电流和电流密度}
    定义单位时间内通过垂直于电流方向单位面积的电量为导体中某点电流密度矢量的$\boldsymbol{j}$大小，$\boldsymbol{j}$的方向与正电荷在该点
    \subsection{欧姆定律的微分形式 焦耳—楞次定律的微分形式}
    \subsection{电源及电源电动势}
    电源电动势在量值

\section{恒定磁场和磁感应强度}
    \subsection{磁性起源于电荷的运动}
    \subsection{磁场 磁感应强度$\boldsymbol{B}$}

\section{毕奥‑萨伐尔定律}
    \subsection{毕奥‑萨伐尔定律}
    在真空中
    \subsection{毕奥‑萨伐尔定律应用举例}
    \subsection{匀速运动电荷的磁场}

\section{磁场中的高斯定理}
    \subsection{磁感应线}
    \subsection{磁通量}
    \subsection{恒定磁场中的高斯定理}

\section{真空中恒定磁场的安培环路定理}
    \subsection{恒定磁场的安培环路定理}
    \subsection{安培环路定理的应用}

\section{磁场对运动电荷和载流导线的作用}
    \subsection{洛伦兹力}
    \subsection{带电粒子在磁场中的运动}
    \subsection{应用电场和磁场控制带电粒子运动的实例}
    \subsection{安培力}

\section{磁力的功}
    \subsection{磁力对运动载流导线的功}
    \subsection{磁力矩对转动载流线圈的功}
    磁力矩对转动载流线圈的功等于电流强度与回路包围面积内的磁通量的增量的乘积。

\section{磁介质中的恒定磁场}
    \subsection{磁介质及其磁化}
    \subsection{磁介质中的高斯定理和安培环路定理}
    \subsection{铁磁质}

\chapter{第八章 电磁场与麦克斯韦电磁场方程组}

\section{电磁感应的基本定律}
    当穿过闭合导体回路的磁通量发生变化时，闭合导体回路中就会出现电流，称之为电磁感应。

    所以

    导体回路中

    应该

    如果，则
    %\[
    %\]
    在$t_1$到$t_2$这段时间为
    %\[
    %\]
    式表明

    总电动势为
    %\[
    %\]

\section{动生电动势}
    把感应电动势
    \subsection{动生电动势的产生原因}
    \subsection{动生电动势的计算}

\section{感生电动势和感生电场}
    \subsection{感生电动势}
    \subsection{感生电场}
    \subsection{感生电动势的计算}
    \subsection{涡电流}
        \subsubsection{热效应}
        \subsubsection{机械效应}
        \subsubsection{趋肤效应}

\section{自感和互感}
    \subsection{自感}
    \subsection{互感}

\section{磁场的能量}

\section{位移电流和全电流定律}
    \subsection{位移电流}
    \subsection{全电流定律}

\section{麦克斯韦电磁场方程组}

\endinput